\subsection{\texorpdfstring{\(C^{r}\) 类函数 ( \(r \neq \omega\) )}{C\^{}\{r\} 类函数 (r ≠ ω)}}

2.3.1. 设 U 是 E 的开子集，f 是从 U 到 F 的映射。对 \(C^{r}\)，按 r 归纳如下定义关系“f 属于 \(r \in N\) 类”：

\begin{enumerate}
\def\labelenumi{\arabic{enumi})}
\item
  \(f\) 属于 \(\mathbf{C}^{0}\) 类，当且仅当 f 连续；
\item
  若 \(r\) 是整数且 r \(\geqslant 1\)，则函数 \(f\) 属于 \(\mathbf{C}^{r}\) 类，当且仅当它在 \(\mathbf{U}\)\(Df\) 的每一点处可微，且从 \(\mathbf{U}\) 到 \(\mathcal{L}(\mathbf{E};\mathbf{F})\) 的导数 Df 属于 \(\mathbf{C}^{r-1}\) 类。
\end{enumerate}

属于 C \(^{r}\) 类的函数也称为 r 次连续可微函数。

若 \(f\)\(\mathbf{C}^{\infty}\) 对每个整数 r 都属于 \(\mathbf{C}^{r}\) 类，则称 f 属于 \(r\) 类（或称无穷次可微）。

若 \(f\) 在 U 中属于 \(\mathbf{C}^{r}\) 类，则对每个整数 \(\mathbf{U}\)，\(f\) 都 \(p\)\(p \leqslant r\) 次可微，且函数 \(\mathbf{D}^{p}f\) 属于 \(\mathbf{C}^{r-p}\) 类。

2.3.2. 从 E 的开子集 U 到 F 的 \(C^{r}\) 类映射构成所有从 U 到 F 的映射空间的向量子空间 \(\mathcal{C}^{r}(U;F)\)。当 \(\mathcal{C}^{s}(U;F)\subset\mathcal{C}^{r}(U;F)\) 时，有 \(s\geqslant r\)。

2.3.3. 函数 \(f\) 在 E 的开子集 U 中属于 \(\mathbf{C}^{1}\) 类的充要条件是它在 U 的每一点处严格可微。

若 E 是赋范空间 \(E_{i}\) 的乘积，则从 E 的开集 V 到 F 的映射 \(f\) 属于 \(C^{r}\) 类，当且仅当对每个整数 m  r，\(f\) 的迭代偏导数 \(D_{i_{1}}\ldots D_{i_{m}}f\)\(m \leqslant r\) 都存在且连续。

2.3.4. 设 G 是赋范空间，U 是 E 的开子集。设 V 是 G 的开子集，\(g \in \mathcal{C}^{r}(U; G)\) 且 \(f \in \mathcal{C}^{r}(V; F)\)。若 \(g(U) \subset V\)，则从 U 到 F 的映射 \(f \circ g\) 属于 \(C^{r}\) 类。

设 \(\mathbf{F}_1\) 和 \(\mathbf{F}_2\) 是两个分离局部凸空间，\(u\) 是从 \(\mathbf{F}_1 \times \mathbf{F}_2\) 到 \(\mathbf{F}\) 的双线性映射，并且关于 \(\mathbf{F}_1\)（分别地，\(\mathbf{F}_2\)）的有界子集族是半连续的（Topological Vector Spaces，第 III 章，第 4 节，第 2 号）。设 \(\mathbf{U}\) 是 \(\mathbf{E}\) 的开子集，且 \(f_i \in \mathcal{C}^r(\mathbf{U};\mathbf{F}_i)\)（\(i = 1,2\)）。则函数 \(u(f_1,f_2)\) 属于 \(\mathcal{C}^r(\mathbf{U};\mathbf{F})\)。若 \(\mathbf{E}\) 有限维，只需假设 \(u\) 满足第 2.1.3 号的条件 (SC)。

2.3.5. 若 E 是赋范空间 \(E_{i}\)（\(1 \leqslant i \leqslant n\)）的乘积，且 f 是从 E 到 F 的连续 n 线性映射，则 f 属于 \(C^{\infty}\) 类，并且当 \(D^{p}f = 0\) 时有 \(p \geqslant n + 1\)。

2.3.6. 假设 E 和 F 是巴拿赫空间。设 f 是在 E 中一点 \(C^{r}\) 的某邻域上定义且取值于 F 的 \(r \geqslant 1\) 类函数（\(x_{0}\)）。设 \(y_{0} = f(x_{0})\)，并假设 \(Df(x_{0})\) 是从 E 到 F 的同构。则 f 诱导出一个从 \(x_{0}\) 的某邻域到 \(y_{0}\) 的某邻域的同胚 g（1.5），且 g 的逆映射在 \(C^{r}\) 的某邻域上属于 \(y_{0}\) 类。

